\documentclass[letterpaper]{article}
\usepackage[submission]{aaai2027}
\usepackage[hyphens]{url}
\usepackage{graphicx}
\urlstyle{rm}
\def\UrlFont{\rm}
\usepackage{natbib}
\frenchspacing
\usepackage{amsmath}
\usepackage{amssymb}
\usepackage{booktabs}

\pdfinfo{
/TemplateVersion (2027.1)
}

\setcounter{secnumdepth}{0}
\newcommand{\sd}[1]{{\footnotesize$\pm$#1}}
\newcommand{\sigmark}[1]{\rlap{\textsuperscript{\scriptsize #1}}}
\newcommand{\tablecaptiongap}{\vspace{0.35em}}

\begin{document}

\title{Course-Knowledge-Guided Reinforcement Learning for Cold-Start Course Recommendation}
\author{Anonymous Submission}
\affiliations{}

\maketitle

\begin{abstract}
MOOC platforms continually publish new courses that must be recommended before they accumulate any interaction history, a persistent cold-start condition. Course recommendation, however, is not only a sparsity problem: a useful new course must also be pedagogically appropriate---resting on mastered prerequisites, matched in difficulty, and non-redundant with prior learning. Generic cold-start recommenders reconstruct item embeddings from content but retain warm-item ranking objectives; when transferred to courses, available educational structure remains outside their original objective. Conversely, course-specific recommenders model learning or prerequisite structure but are not designed around zero-interaction target courses. We propose Course-Knowledge-Guided Reinforcement Learning (CKG-RL), which casts cold-start course ranking as a sequential candidate-selection policy trained to maximize an educational reward rather than to fit past clicks, placing concept overlap, prerequisites, difficulty, and redundancy inside both candidate construction and the reward; it further anchors course representations in content, masks cold-course ID evidence, and refines course states through retrieval-based learner simulation. We evaluate under strict course-cold conditions---zero-interaction target courses competing against the entire catalog, scored by course-macro Recall/NDCG so that each new course, rather than each interaction of a popular course, is the unit of success. On MOOCCube, Junyi, and COCO, CKG-RL consistently improves early-rank exposure of genuinely new courses over strong content-transfer and graph-reconstruction baselines, with gains that are statistically significant under Holm-corrected paired tests in nearly all early-rank comparisons. MOOCCube ablations attribute the gains to educational rewards, knowledge-guided sampling, and prerequisite supervision.
\end{abstract}

\section{Introduction}
MOOC platforms continually add courses, course versions, and learning paths. Recommenders therefore must rank courses before interaction histories are available, a persistent cold-start condition~\cite{schein2002coldstart,gantner2010attribute,lin2024tdro,zhang2025coldstartllm}. In education, cold-start ranking is not only a sparsity problem. A useful course should be topically relevant, prerequisite-safe, difficulty-appropriate, and non-redundant with prior learning. Public MOOC resources expose such course-side signals through metadata, concepts, and prerequisite relations~\cite{yu2020mooccube,wang2021mooccubex}.

Existing work addresses the two requirements separately. Generic cold-start methods reconstruct item representations from content or side attributes~\cite{schein2002coldstart,gantner2010attribute,volkovs2017dropoutnet,kim2024cgrc}, whereas course-specific methods model learning decisions or prerequisites~\cite{zhang2019hrlcourse,sun2024pcgnn}. We therefore use CGRC as a generic cold-start transfer baseline, not a course recommender, and PCGNN as the course-specific counterpart. Figure~\ref{fig:method-motivation} diagnoses their complementary limitations on held-out MOOCCube tests. PCGNN retains prerequisite and difficulty mismatch under zero interaction, while 46\% of the audited CGRC cold-course cases have NDCG@10 no greater than 0.10 and cold courses occupy only 24.6\% of its Top-10 lists. These observations motivate course-knowledge constraints for structural fit and cold-course representation learning for effective full-catalog ranking.

To measure this honestly, we study \emph{strict course-cold MOOC recommendation}: target courses have zero training interactions, compete against the full catalog, and are evaluated by course-macro Recall/NDCG, so the evaluation unit is the cold course itself rather than a cold-course interaction drawn from a popular course. This distinction matters because common offline protocols can hide the failure above: sampled candidate sets remove the warm courses that dominate the real catalog~\cite{krichene2020sampled}, weak protocols and baselines can overstate progress~\cite{dacrema2019progress}, and interaction-weighted averages can be driven by a few high-volume cold courses~\cite{takehi2025diversification}. The benchmark thus forces the model to answer two questions at once: whether a new course obtains early exposure at all, and whether the exposed course is educationally plausible for the learner.

\begin{figure}[t]
\centering
\includegraphics[width=\columnwidth]{figures/mooccube_method_motivation.pdf}
\caption{\textbf{Why existing methods fall short.} On held-out MOOCCube tests, PCGNN retains structural mismatch, while 46\% of CGRC cold-course cases have NDCG@10 no greater than 0.10 and cold courses occupy 24.6\% of its Top-10 lists. CKG-RL targets these gaps through course-knowledge constraints and cold-course representation learning; Figure~\ref{fig:p1-topk-motivation} evaluates the resulting improvements and trade-offs.}
\label{fig:method-motivation}
\end{figure}

CKG-RL responds with content anchoring and cold-ID masking for the evidence boundary, learner-state simulation for cold representations, and knowledge-guided sampling, rewards, and prerequisite supervision for ranking constraints.

Across MOOCCube, Junyi, and COCO, CKG-RL has higher @10 course-macro means than the strongest non-CKG baseline under strict course-cold full-ranking evaluation. We restrict this conclusion to early-rank exposure for zero-interaction courses. Some baselines remain competitive on broader or hot-course diagnostics.

This paper makes three contributions:
\begin{itemize}
    \item We define a strict course-cold MOOC benchmark with full-catalog course-macro Recall/NDCG as the primary target.
    \item We introduce CKG-RL, which combines content anchoring, forced-cold masking, retrieval-based simulation, and educational rewards.
    \item We evaluate three datasets against collaborative, graph, content-transfer, and cold-start baselines; ablations isolate the course-knowledge components.
\end{itemize}

\section{Related Work}
Educational recommendation must match resources with pedagogical context and learning paths~\cite{dasilva2023educationalrs,rahayu2023learningpath}. MOOCCube and MOOCCubeX provide interactions, course metadata, concepts, teachers, schools, and prerequisite relations~\cite{yu2020mooccube,wang2021mooccubex}. Prior MOOC recommenders have used sequential decisions and explainable path reasoning~\cite{zhang2019hrlcourse,dessi2024findingpaths}. For course-cold targets, these course-side signals remain available even though target-course behavior is absent. The available evidence is therefore course-side. We treat concepts, prerequisites, difficulty, and redundancy as evidence for cold-course ranking.

Cold-start recommendation maps side attributes and content into collaborative spaces when observations are sparse~\cite{schein2002coldstart,gantner2010attribute,liu2024once,wu2024llmrec,zhang2025coldstartllm}. Interaction-only and graph recommenders such as BPR, NGCF, and LightGCN~\cite{rendle2009bpr,wang2019ngcf,he2020lightgcn} are strong when enough edges are observed, but course cold-start requires ranking targets whose IDs provide little supervised signal. Generic cold-start methods transfer warm-item information through dropout, contrastive alignment, distillation, graph reconstruction, LLM augmentation or simulation, OOV refinement, structure completion, meta-learning, or exploration~\cite{volkovs2017dropoutnet,zhou2023ccfcrec,huang2023aldi,kim2024cgrc,wang2024llmaugmenter,huang2025coldllm,zhao2024usim,zhao2025fsgnn,luo2025pam,jiang2025itemcentric}. Our question is whether this transfer view suffices after adaptation to full-catalog course ranking with educational constraints.

Offline recommender evaluation is sensitive to metric choice, candidate construction, and reproducibility practice~\cite{sun2020rigorous,hidasi2023flaws,bauer2024landscape}. Sampled negatives can change comparisons relative to full ranking~\cite{krichene2020sampled,li2024itemsampling}, while distribution-aware evaluation exposes variation across poorly served groups or courses~\cite{ekstrand2024distributional}. We apply this concern to course-cold recommendation: the unit at risk is a held-out course, and the failure is no useful exposure under full-catalog ranking~\cite{takehi2025diversification}. Reinforcement-learning recommenders model sequential decisions~\cite{shani2005mdp,zhao2018deepreinforcement,chen2023rlrecsurvey}. CKG-RL uses a narrower simulator for learner-course matching under course-cold evidence constraints.

\section{Problem Formulation}

\subsection{Preliminaries}
Let \(\mathcal{U}\) denote the set of learners and \(\mathcal{C}\) denote the set of courses. The static training interactions satisfy \(\mathcal{D}_{\mathrm{train}}\subseteq\mathcal{U}\times\mathcal{C}\); an observed pair \((u,c)\) records that learner \(u\in\mathcal{U}\) interacted with course \(c\in\mathcal{C}\). Given \(\mathcal{D}_{\mathrm{train}}\) and course-side information, a recommender produces, for each learner \(u\), a ranked list \(\Pi_u\) over a candidate course set (specified in Eq.~\eqref{eq:candidate-set}).

Each course \(c\) has a content vector \(\mathbf{x}_c\) and educational signals. We use concept overlap, prerequisite links, training popularity, and history similarity as proxies for topical continuity, readiness, difficulty, and redundancy, consistent with pedagogical-fit views of educational recommendation~\cite{dasilva2023educationalrs,rahayu2023learningpath}. These signals are available even when \(c\) has no training interactions. They are not complete learner-state measurements.

\subsection{Strict Course-Cold Protocol}
We study a course-cold held-out setting. The course set is partitioned into warm courses \(\mathcal{C}_{\mathrm{warm}}\) and cold test courses \(\mathcal{C}_{\mathrm{cold}}\). A cold test course satisfies
\begin{equation}
|\{(u',c')\in \mathcal{D}_{\mathrm{train}}: c'=c\}|=0.
\end{equation}
Thus the model cannot use interaction-trained evidence for the target course ID. Concretely, the target-course ID embedding is replaced by a masked vector \(\bar{\mathbf{v}}_c\) (defined in Eq.~\eqref{eq:cold-mask}), so its representation must come from content and educational knowledge.

For a test pair \((u,c^+)\) with \(c^+\in\mathcal{C}_{\mathrm{cold}}\), let \(\mathcal{H}^{\mathrm{train}}_u\) be the courses observed in \(u\)'s training history. The available candidate set is the full catalog after this history is removed:
\begin{equation}
\mathcal{A}_u=\mathcal{C}\setminus \mathcal{H}^{\mathrm{train}}_u .
\label{eq:candidate-set}
\end{equation}
The model scores every course in \(\mathcal{A}_u\), avoiding sampled-candidate artifacts~\cite{krichene2020sampled,li2024itemsampling} and weak-protocol effects that can distort recommender comparisons~\cite{dacrema2019progress,hidasi2023flaws}. The split is course-based rather than user-based: test learners have training histories, while target courses have none.

\subsection{Standard Interaction Average}
Let \(m_K(u,c^+)\) denote a top-\(K\) utility for a test pair, instantiated as Recall@\(K\) or NDCG@\(K\). Let \(\kappa_{u,c^+}\) be the rank position of target course \(c^+\) in \(\Pi_u\). Recall@\(K\) equals one when \(\kappa_{u,c^+}\le K\) and zero otherwise, and NDCG@\(K\) is
\begin{equation}
\operatorname{NDCG@K}(u,c^+)=
\begin{cases}
1/\log_2(\kappa_{u,c^+}+1), & \kappa_{u,c^+}\le K,\\
0, & \kappa_{u,c^+}>K.
\end{cases}
\label{eq:topk-utility}
\end{equation}
Since each test case contains one held-out positive course, these metrics measure whether the target cold course is surfaced early in the full ranking.

Let \(\mathcal{D}_{\mathrm{test}}\) be the held-out interaction set and \(\mathcal{D}^{\mathrm{cold}}_{\mathrm{test}}:=\{(u,c^+)\in\mathcal{D}_{\mathrm{test}}:c^+\in\mathcal{C}_{\mathrm{cold}}\}\). A common evaluation choice is to average the utility over these cold-course test interactions:
\begin{equation}
\operatorname{Micro}(K)=
\frac{1}{|\mathcal{D}^{\mathrm{cold}}_{\mathrm{test}}|}
\sum_{(u,c^+)\in \mathcal{D}^{\mathrm{cold}}_{\mathrm{test}}}
m_K(u,c^+).
\end{equation}
This measures interaction-level utility, but it does not show whether each cold course receives exposure.

\subsection{Cold-Course Risk Metric}
Strict course-cold recommendation has a course-level failure mode: high-volume cold courses can dominate the average while other cold courses remain unseen. To expose this tail risk and distributional variation~\cite{ekstrand2024distributional,takehi2025diversification}, we aggregate first within each cold course. Let \(\mathcal{D}_{\mathrm{test}}(c):=\{(u,c^+)\in\mathcal{D}^{\mathrm{cold}}_{\mathrm{test}}:c^+=c\}\) be the test cases whose held-out target course is \(c\), and let \(\mathcal{C}^{\mathrm{eval}}_{\mathrm{cold}}:=\{c\in\mathcal{C}_{\mathrm{cold}}:\mathcal{D}_{\mathrm{test}}(c)\neq\emptyset\}\) be the cold courses that appear as a target at test time. For \(c\in\mathcal{C}^{\mathrm{eval}}_{\mathrm{cold}}\), the scalar course-level quality at cutoff \(K\) is
\begin{equation}
q_c(K)=
\frac{1}{|\mathcal{D}_{\mathrm{test}}(c)|}
\sum_{(u,c)\in \mathcal{D}_{\mathrm{test}}(c)}
m_K(u,c),
\end{equation}
where each case in \(\mathcal{D}_{\mathrm{test}}(c)\) has \(c\) as its held-out positive, so \(m_K(u,c)\) is the top-\(K\) utility defined around Eq.~\eqref{eq:topk-utility}. The course-macro metric is then
\begin{equation}
\operatorname{Macro}(K)=
\frac{1}{|\mathcal{C}^{\mathrm{eval}}_{\mathrm{cold}}|}
\sum_{c\in\mathcal{C}^{\mathrm{eval}}_{\mathrm{cold}}} q_c(K).
\end{equation}
Micro averaging asks whether a random cold-course interaction is served well. Course-macro averaging asks whether a random cold course is served well. We use course-macro as the primary metric because the unit of failure is a zero-interaction course receiving no useful full-catalog rank.

\subsection{Main Objective}
The objective is to learn a logit \(\ell(u,c)\), realized as the learner--course match of Eq.~\eqref{eq:ranking-logit}, that ranks \(c^+\in\mathcal{C}_{\mathrm{cold}}\) highly within \(\mathcal{A}_u\), using learner history and course-side evidence but no interaction-derived target-course signal. Hot-course metrics are diagnostics for the trade-off between zero-interaction generalization and warm-course ranking.

\section{Method}
\subsection{Overview}
CKG-RL makes the evidence boundary of strict course-cold ranking explicit. For each learner-course pair \((u,c)\), the model builds a learner representation \(\mathbf{z}_u\) from the learner history and an initial course embedding \(\mathbf{e}_c\) from course-side evidence. A simulator starts from \(\mathbf{e}_c\), refines it through \(T\) learner-matching steps, and returns the final course embedding \(\mathbf{z}_{c,T}:=\mathbf{h}_{T,c}\). The ranking logit \(\ell(u,c)\) then matches \(\mathbf{z}_{c,T}\) with the target learner representation \(\mathbf{z}_u\) (Eq.~\eqref{eq:ranking-logit}).

Three design constraints connect the architecture to the two-part motivation. Content anchoring and forced masking enforce the zero-interaction evidence boundary. Simulation exposes cold representations to plausible learner states under full-catalog competition. Course knowledge then filters matches that are semantically related but prerequisite-unsafe, too difficult, or redundant. Figure~\ref{fig:ckg-rl-framework} summarizes the training and evaluation flow.

\begin{figure*}[t]
\centering
\includegraphics[width=0.98\textwidth]{figures/ckg_rl_framework_topconf.pdf}
\caption{Overview of CKG-RL for strict course-cold ranking. Cold-course ID evidence is masked, course states are initialized from content, refined through knowledge-guided retrieval, and ranked against the full catalog.}
\label{fig:ckg-rl-framework}
\end{figure*}

\subsection{Content-Anchored Course Representation}
The content branch provides the evidence available for a course with zero training interactions, a role emphasized in content-based and LLM-assisted recommendation~\cite{liu2024once,wu2024llmrec}. Each learner \(u\) has a trainable ID embedding \(\mathbf{p}_u\), projected into the recommendation space:
\begin{equation}
\mathbf{z}_u=f_u(\mathbf{p}_u).
\end{equation}
Each course \(c\) has a trainable ID embedding \(\mathbf{v}_c\) and a frozen 768-dimensional content embedding \(\mathbf{x}_c\). A content projection network maps \(\mathbf{x}_c\) into the same latent space:
\begin{equation}
\mathbf{c}_c=f_c(\mathbf{x}_c).
\end{equation}
The initial course embedding is a gated fusion of the projected content vector and a possibly masked ID vector:
\begin{equation}
\mathbf{g}_c=\sigma\left(\mathbf{W}_g[\bar{\mathbf{v}}_c;\mathbf{c}_c]\right),
\quad
\mathbf{e}_c=\mathbf{g}_c\odot \bar{\mathbf{v}}_c+(1-\mathbf{g}_c)\odot \mathbf{c}_c.
\end{equation}
Here \(f_u\) and \(f_c\) are learned projection networks, \(\mathbf{W}_g\) is a learned gate matrix, \(\sigma\) is the sigmoid function, and \(\odot\) denotes elementwise multiplication. \(\bar{\mathbf{v}}_c\) is the ID branch after masking and ID dropout. The gate lets warm courses use collaborative evidence when available, while the content branch defines \(\mathbf{e}_c\) when ID evidence is unavailable.

\subsection{Forced-Cold ID Masking}
Forced-cold masking enforces the protocol inside the model. Let \(p_c\) be the training-set popularity of course \(c\). Under the strict course-cold protocol, a course is cold when \(p_c<1\), which is equivalent to having zero training interactions. For cold examples, CKG-RL removes the course ID evidence:
\begin{equation}
\bar{\mathbf{v}}_c=(1-\xi_c)\mathbf{v}_c,
\quad
\xi_c=\mathbb{I}[p_c<1].
\label{eq:cold-mask}
\end{equation}
Here \(\mathbb{I}[\cdot]\) is the indicator function. During training, ordinary ID dropout is additionally applied with probability \(0.35\). At cold-course evaluation time, the target course vector is also built with the ID branch masked. This keeps the representation used at inference aligned with the evidence available in the strict course-cold protocol.

\subsection{Retrieval-Based User Simulation}
Full-catalog course-cold ranking requires a course representation that can compete for plausible learners, not only a pointwise content score. Like recent augmentation or simulation approaches~\cite{wang2024llmaugmenter,huang2025coldllm}, CKG-RL compensates for missing course interactions while keeping the simulator bounded. It initializes \(\mathbf{h}_{0,c}=\mathbf{e}_c\). At step \(t\), it retrieves candidate learners from a normalized user bank by cosine similarity to \(\mathbf{h}_{t,c}\). Let \(\mathcal{N}_M(\mathbf{h}_{t,c})\) be the top-\(M\) retrieved learners. CKG-RL samples \(N\) candidates from this set with a temperature-smoothed distribution \(q_t\) and an exploration floor:
\begin{equation}
\begin{aligned}
q_t(u'\mid \mathbf{h}_{t,c})=&
(1-\epsilon)
\frac{\exp(\mathbf{h}_{t,c}^{\top}\mathbf{z}_{u'}/\tau_s)}
{Z_t}
+\frac{\epsilon}{M},\\
Z_t=&
\sum_{v\in\mathcal{N}_M(\mathbf{h}_{t,c})}
\exp(\mathbf{h}_{t,c}^{\top}\mathbf{z}_{v}/\tau_s),\\
&u'\in\mathcal{N}_M(\mathbf{h}_{t,c}).
\end{aligned}
\end{equation}
Here \(\tau_s>0\) is the sampling temperature and \(\epsilon\in[0,1]\) is the exploration mass. Let \(w_C,w_P,w_D,w_R\geq0\) weight the concept bonus, prerequisite gap, difficulty gap, and redundancy penalty, respectively. The sampled candidates are then reordered by a course-fit score:
\begin{equation}
\begin{aligned}
b(u',c)=&
w_C C(u',c)-w_P P(u',c)\\
&-w_D D(u',c)-w_R R(u',c),
\end{aligned}
\end{equation}
\begin{equation}
s_{sample}(u',c)=
\cos(\mathbf{h}_{t,c},\mathbf{z}_{u'})
+\beta\operatorname{scale}_t(b(u',c)).
\end{equation}
Here \(\operatorname{scale}_t(\cdot)\) maps a scalar fit score to \([0,1]\) over the step-\(t\) sampled candidates, \(\beta\geq0\) weights the knowledge-fit bias relative to latent similarity, and \(C\), \(P\), \(D\), and \(R\) are the reward terms defined below. This knowledge-guided sampling bias is applied to both cold and hot examples in the final run, so retrieval is filtered by educational plausibility as well as latent similarity.

An actor-critic action distribution \(\pi_t(\cdot\mid\mathbf{s}_t)\) selects a learner \(a_t\) from the sampled candidates. The simulator state concatenates the current course representation with the normalized step index, \(\mathbf{s}_t=[\mathbf{h}_{t,c};\,t/(T-1)]\), so the policy is aware of its position within the \(T\)-step horizon. CKG-RL then updates the course state by taking a step of size \(\eta>0\) toward a mixture of the selected learner and the initial course embedding:
\begin{equation}
\mathbf{h}_{t+1,c}=\mathbf{h}_{t,c}+\eta
\nabla_{\mathbf{h}_{t,c}}
\left[
(1-\alpha_t)\mathbf{h}_{t,c}^\top\mathbf{z}_{a_t}
+\alpha_t\mathbf{h}_{t,c}^\top\mathbf{e}_c
\right].
\end{equation}
The target coefficient is adaptive:
\begin{align}
\alpha_t=
\operatorname{clip}_{[\alpha_{\min},\alpha_{\max}]}\Big(
&\alpha_{cold}\xi_c+\alpha_{hot}(1-\xi_c) \nonumber\\
&+\alpha_{step}\frac{t}{T-1}
-\alpha_{ent}\frac{H(\pi_t)}{\log N}
\Big).
\end{align}
where \(\operatorname{clip}_{[a,b]}(x)=\min\{b,\max\{a,x\}\}\) and \(H(\pi_t)\) is the entropy of the action distribution. The base coefficients \(\alpha_{cold}\) and \(\alpha_{hot}\) set the target-anchoring strength for cold and hot courses, \(\alpha_{step}\ge0\) increases target guidance as the step index grows, \(\alpha_{ent}\ge0\) attenuates anchoring when the policy is uncertain, and \([\alpha_{\min},\alpha_{\max}]\) bounds the result.
Thus cold courses receive stronger target anchoring, later updates within the simulator horizon increase target guidance, and high-entropy actions reduce overconfident target drift. The simulator does not reconstruct a real learner trajectory. It is a training-time mechanism that exposes the cold-course representation to a sequence of plausible learner matches under full-catalog competition; the sequential structure lets the adaptive coefficient \(\alpha_t\) tighten target anchoring as evidence accumulates over the horizon, which a single pointwise match cannot express.

\subsection{Course-Knowledge Reward Design}
The reward terms convert educational constraints into simulator signals, reflecting the learning-resource literature's emphasis on context, sequencing, and learner-resource fit~\cite{dasilva2023educationalrs,rahayu2023learningpath}. Let \(\mathbf{y}_u\in\{0,1\}^{|\mathcal{C}|}\) be the multi-hot vector of courses seen by learner \(u\), with \(n_u=\sum_{c'\in\mathcal{C}}y_{u,c'}\). For target course \(c\), \(\mathbf{a}^{\mathrm{pre}}_c\in\{0,1\}^{|\mathcal{C}|}\) marks its prerequisite courses and \(\mathbf{a}^{\mathrm{con}}_c\in[0,1]^{|\mathcal{C}|}\) gives its concept overlap with each course. The prerequisite gap estimates missing prerequisite evidence:
\begin{equation}
P(u,c)=
\left[
1-\frac{\mathbf{y}_u^\top \mathbf{a}^{\mathrm{pre}}_c}
{\max(1,\|\mathbf{a}^{\mathrm{pre}}_c\|_1)}
\right]_0^1.
\end{equation}
Write \([x]_0^1:=\min\{1,\max\{0,x\}\}\). A prerequisite-safe gate is \(S_{\mathrm{pre}}(u,c)=\mathbb{I}[P(u,c)\le\delta_{\mathrm{pre}}]\), where \(\delta_{\mathrm{pre}}\) is the safety threshold. The concept match is
\begin{equation}
\chi(u,c)=
\left[
\frac{\mathbf{y}_u^\top \mathbf{a}^{\mathrm{con}}_c}{\max(1,n_u)}
\right]_0^1.
\end{equation}
Let \(\mu_C\) and \(\mu_R\) be the lower and redundancy concept-overlap thresholds, with \(0\le\mu_C<\mu_R<1\), and let \(\gamma_R\in[0,1]\) be the redundancy attenuation.
Course difficulty is estimated from inverse training popularity,
\begin{equation}
d_c=1-\frac{\log(1+p_c)}
{\max_{c'\in\mathcal{C}} \log(1+p_{c'})}.
\end{equation}
with learner readiness \(r_u=[n_u/n_{\mathrm{sat}}]_0^1\), where \(n_{\mathrm{sat}}\) is the history length at which readiness saturates, and gap \(D(u,c)=\max(0,d_c-r_u)\). When concept overlap is too high, the redundancy penalty is
\begin{equation}
R(u,c)=
\left[
\frac{\chi(u,c)-\mu_R}{1-\mu_R}
\right]_0^1.
\end{equation}
The bounded concept bonus is
\begin{equation}
C(u,c)=
\left[
\frac{\chi(u,c)-\mu_C}{\mu_R-\mu_C}
\right]_0^1
\cdot S_{\mathrm{pre}}(u,c)\cdot (1-\gamma_R R(u,c)).
\end{equation}

The step reward combines target preservation, progress, duplicate-action control, and the four course-knowledge terms:
\begin{align}
r_t=&-w_T\|\mathbf{h}_{t+1,c}-\mathbf{e}_c\|_2^2
+w_G G_t-w_{dup}\operatorname{Dup}_t \nonumber\\
&+w_C C(a_t,c)-w_P P(a_t,c) \nonumber\\
&-w_D D(a_t,c)-w_R R(a_t,c),
\end{align}
where \(w_T,w_G,w_{\mathrm{dup}}\) weight target preservation, progress, and repeated-action control, respectively; \(\operatorname{Dup}_t\in\{0,1\}\) indicates a repeated action; and \(G_t=[\|\mathbf{h}_{t,c}-\mathbf{e}_c\|_2^2-\|\mathbf{h}_{t+1,c}-\mathbf{e}_c\|_2^2]_0^1\).

\subsection{Optimization}
Optimization combines policy learning with ranking supervision. Let \(V_\phi(\mathbf{s}_t)\) be the critic, \(\gamma\in[0,1]\) the discount factor, and \(\lambda\in[0,1]\) the GAE decay. The actor-critic policy is optimized with PPO and generalized advantage estimation~\cite{schulman2017ppo,schulman2015gae}:
\begin{equation}
\delta_t=r_t+\gamma V_\phi(\mathbf{s}_{t+1})-V_\phi(\mathbf{s}_t),
\quad
\hat{A}_t=\sum_{j\geq 0}(\gamma\lambda)^j\delta_{t+j}.
\end{equation}
With \(\rho_t(\theta)=\pi_\theta(a_t\mid\mathbf{s}_t)/\pi_{\theta_{\mathrm{old}}}(a_t\mid\mathbf{s}_t)\), where \(\pi_\theta\) is the actor policy, \(\pi_{\theta_{\mathrm{old}}}=\pi_t\) is the rollout policy, and \(\epsilon_p>0\) is the PPO clipping radius, the clipped actor objective is
\begin{equation}
\mathcal{L}_{actor}=
-\mathbb{E}_t\left[
\min(\rho_t\hat{A}_t,
\operatorname{clip}(\rho_t,1-\epsilon_p,1+\epsilon_p)\hat{A}_t)
\right].
\end{equation}
After the \(T\)-step simulator horizon, we denote the final course embedding as \(\mathbf{z}_{c,T}:=\mathbf{h}_{T,c}\). The ranking logit is
\begin{equation}
\ell(u,c)=\frac{\mathbf{z}_u^\top\mathbf{z}_{c,T}}
{\tau\|\mathbf{z}_u\|_2\|\mathbf{z}_{c,T}\|_2},
\label{eq:ranking-logit}
\end{equation}
where \(\tau>0\) is the ranking temperature. The main ranking loss is a cross-entropy contrastive loss over the positive course and mixed hard/random in-batch negatives. Same-course and known-positive negatives are masked to reduce false negatives, a known concern in negative-sampling practice~\cite{ma2026negative}. The model also applies an auxiliary InfoNCE loss between the course ID embedding and the content embedding, plus a prerequisite auxiliary margin loss for high-violation negatives. These terms keep the ranking objective tied to observed learner-course pairs while PPO shapes the simulation policy. The final objective is
\begin{equation}
\mathcal{L}=
\mathcal{L}_{rank}
+\lambda_{aux}\mathcal{L}_{aux}
+\mathcal{L}_{ppo}
+\lambda_{pre}\mathcal{L}_{pre}.
\end{equation}
Here \(\mathcal{L}_{rank}\), \(\mathcal{L}_{aux}\), \(\mathcal{L}_{ppo}\), and \(\mathcal{L}_{pre}\) are the ranking, content-alignment, PPO, and prerequisite-margin losses, and \(\lambda_{aux},\lambda_{pre}\geq0\) weight the two auxiliary losses.

\subsection{Implementation Details}
The implementation uses 128-dimensional learner and course embeddings, a 256-dimensional hidden layer, and 768-dimensional course content features. The MOOCCube/MOOCCubeX and COCO content features are BERT [CLS] representations~\cite{devlin2019bert}, using Chinese and English encoders respectively; Junyi is treated as an educational course generalization setting with processed course-content features. The simulator uses horizon \(T=5\) with update rate \(\eta=0.3\). At each step, it retrieves \(M=256\) users, samples \(N=20\) candidates with \(\tau_s=0.20\) and \(\epsilon=0.10\), and reorders candidates with knowledge-sampling weight \(\beta=0.20\). The educational reward weights are \(0.08\) for prerequisite penalty, \(0.04\) for concept bonus, \(0.03\) for difficulty penalty, and \(0.02\) for redundancy penalty. We use \(n_{\mathrm{sat}}=5\), \(\delta_{\mathrm{pre}}=0.20\), \(\mu_C=0.12\), \(\mu_R=0.70\), and \(\gamma_R=1.0\). The target-anchor coefficients are \(\alpha_{cold}=0.35\), \(\alpha_{hot}=0.60\), \(\alpha_{step}=0.20\), \(\alpha_{ent}=0.20\), and \([\alpha_{\min},\alpha_{\max}]=[0.15,0.85]\). PPO uses two epochs, \(\gamma=0.90\), \(\lambda=0.95\), clipping radius \(\epsilon_p=0.20\), value clipping \(0.20\), and advantage normalization. The final strict course-cold runs train for 60 epochs with batch size 2048 and select checkpoints using cold course-macro validation behavior. The reported configuration uses only the components described in this section.

\section{Experiments}
\subsection{Dataset and Protocol}
The experiments use three educational recommendation settings, summarized in Table~\ref{tab:dataset-statistics}. MOOCCube is the course-knowledge analysis setting: it provides learner-course interactions, and MOOCCubeX supplies course-concept and prerequisite evidence~\cite{yu2020mooccube,wang2021mooccubex}. Junyi is included as a third educational course recommendation setting under the same strict course-cold evaluator~\cite{chen2020junyi}, but it is not used for the course-relation ablations. COCO is the large-catalog main-table setting and provides a public semantic collection of online courses with a preprocessed knowledge graph over course metadata~\cite{dessi2018coco}.

\begin{table}[t]
\centering
{
\small
\setlength{\tabcolsep}{4pt}
\begin{tabular}{lrrrr}
\toprule
\textbf{Datasets} & \textbf{\#Users} & \textbf{\#Courses} & \textbf{\#Inter.} & \textbf{Density} \\
\midrule
\textbf{MOOCCube} & 199,199 & 698 & 682,753 & 0.4910\% \\
\textbf{Junyi} & 156,244 & 722 & 2,198,427 & 1.9500\% \\
\textbf{COCO} & 24,036 & 8,196 & 378,469 & 0.1921\% \\
\bottomrule
\end{tabular}
}
\caption{Statistics of strict course-cold evaluation datasets.}
\label{tab:dataset-statistics}
\end{table}

The strict course-cold balanced protocol uses three random seeds, 2025, 2026, and 2027. Warm interactions are split 80/10/10 before cold-course reinsertion. Eligible cold courses are assigned to 20 popularity-balanced folds: one validation fold and two test folds, giving validation/test cold course ratios of 0.05/0.10. All interactions of validation and test cold courses are removed from training, while test users must retain training history, isolating course rather than user cold-start. The final course-macro evaluation contains about 68 MOOCCube cold courses, 71 Junyi cold courses, and 820 COCO cold courses per seed. After train-history masking, COCO uses about 6,048 warm candidate courses per seed.

The split is a controlled course-cold benchmark, not a chronological launch-time study. Primary metrics are full-ranking course-macro Recall@\(K\) and NDCG@\(K\), \(K \in \{5,10,20\}\). Hot course-macro metrics are auxiliary diagnostics. This separates zero-interaction generalization, full-catalog competition, and whether gains cover cold courses rather than concentrating on high-volume ones~\cite{bauer2024landscape,ekstrand2024distributional}.

\subsection{Research Questions}
The experiments answer four questions. \textbf{RQ1:} What strict course-cold full-ranking course-macro scores does CKG-RL obtain relative to the baselines? \textbf{RQ2:} Which exposure and pedagogical-risk patterns motivate the method components? \textbf{RQ3:} Which course-knowledge components contribute, and do their list-level effects match their intended roles? \textbf{RQ4:} How stable are the gains under design variations, and what cost does CKG-RL introduce?

\subsection{Baselines}
To keep the comparison protocol-controlled, all reported cells use the same course-cold splits, full-catalog candidate set, train-history masking, course-macro evaluator, and validation cold NDCG@10 checkpoint rule~\cite{sun2020rigorous,hidasi2023flaws}. We compare against several capability classes. \textbf{Interaction-only} baselines include Popularity, BPR~\cite{rendle2009bpr}, and LightGCN~\cite{he2020lightgcn}. They test whether learner-side history or graph propagation can rank target courses whose own training edges are removed. \textbf{Side-feature fallback} uses DropoutNet~\cite{volkovs2017dropoutnet}. \textbf{Content-transfer} uses CCFCRec~\cite{zhou2023ccfcrec}, ALDI~\cite{huang2023aldi}, and SEMCo~\cite{meehan2026semco}. \textbf{Knowledge-graph and graph-reconstruction} baselines include KGRec~\cite{yang2023knowledge}, generic cold-start CGRC~\cite{kim2024cgrc}, and the adapted course-specific PCGNN~\cite{sun2024pcgnn}. CGRC is closest in cold-start transfer mechanism; PCGNN supplies the course-specific prerequisite counterpart. DropoutNet, ALDI, and SEMCo retain their released objectives under the common evaluator. KGRec follows its original objective, CCFCRec uses a public-source static-HIN wrapper, and CGRC is reimplemented because no official code was available.

\subsection{Theoretical Complexity}
Let \(E\) be the number of training interactions, \(I\) the catalog size, \(Q\) the number of evaluation cases, and \(d\) the embedding dimension. CKG-RL adds a bounded simulator term \(O(E T(Md+Nd+H))\) during training, where \(T=5\), \(M=256\), and \(N=20\) in the reported runs. After course representations are cached, full-ranking inference uses the same dominant dot-product term \(O(QId)\) as embedding baselines. The supplement reports seed-level CIs, per-course gain diagnostics, and CGRC cost/timing coverage.

\subsection{RQ1: Strict Course-Cold Accuracy}
Table~\ref{tab:main-cold-results} reports the primary cross-dataset comparison for Recall/NDCG at the early-rank cutoffs @5 and @10. Full @5/@10/@20 results, seed-level CIs, and per-course gain diagnostics are in the supplement. Across all reported cells under the same full-catalog course-macro evaluator and matched three-seed protocol, CKG-RL has the highest mean in every @5/@10 column on MOOCCube, Junyi, and COCO. The COCO interpretation is restricted to early-rank exposure because CCFCRec remains stronger at Recall@20 in the supplementary table.

\begin{table*}[t]
\centering
{
\small
\setlength{\tabcolsep}{1.35pt}
\renewcommand{\arraystretch}{0.98}
\begin{tabular}{lcccccccccccc}
\toprule
& \multicolumn{4}{c}{MOOCCube} & \multicolumn{4}{c}{Junyi} & \multicolumn{4}{c}{COCO} \\
\cmidrule(lr){2-5}\cmidrule(lr){6-9}\cmidrule(lr){10-13}
Method & R@5 & R@10 & N@5 & N@10 & R@5 & R@10 & N@5 & N@10 & R@5 & R@10 & N@5 & N@10 \\
\midrule
Popularity & 0.0000 & 0.0000 & 0.0000 & 0.0000 & 0.0000 & 0.0000 & 0.0000 & 0.0000 & 0.0000 & 0.0000 & 0.0000 & 0.0000 \\
BPR & 0.0113 & 0.0249 & 0.0068 & 0.0111 & 0.0298 & 0.0578 & 0.0180 & 0.0270 & 0.0000 & 0.0000 & 0.0000 & 0.0000 \\
DropoutNet & 0.1166 & 0.1658 & 0.0791 & 0.0950 & 0.0540 & 0.0928 & 0.0336 & 0.0460 & 0.0391 & 0.0650 & 0.0243 & 0.0326 \\
LightGCN & 0.0833 & 0.1173 & 0.0593 & 0.0702 & 0.0187 & 0.0420 & 0.0108 & 0.0182 & 0.0067 & 0.0132 & 0.0040 & 0.0060 \\
CCFCRec & 0.1476 & 0.2033 & 0.1101 & 0.1281 & 0.0942 & 0.1462 & 0.0619 & 0.0786 & \underline{0.0393} & \underline{0.0656} & \underline{0.0248} & \underline{0.0332} \\
ALDI & 0.1733 & 0.2258 & 0.1213 & 0.1382 & \underline{0.1331} & \underline{0.2102} & \underline{0.0859} & \underline{0.1106} & 0.0202 & 0.0383 & 0.0121 & 0.0179 \\
KGRec & 0.0974 & 0.1305 & 0.0764 & 0.0871 & 0.0245 & 0.0436 & 0.0149 & 0.0210 & 0.0070 & 0.0126 & 0.0042 & 0.0060 \\
CGRC & \underline{0.2121} & \underline{0.2589} & \underline{0.1695} & \underline{0.1845} & 0.1269 & 0.1915 & 0.0833 & 0.1041 & 0.0250 & 0.0377 & 0.0168 & 0.0209 \\
PCGNN & 0.0241 & 0.0520 & 0.0145 & 0.0234 & 0.0453 & 0.0897 & 0.0261 & 0.0402 & 0.0135 & 0.0254 & 0.0081 & 0.0118 \\
USIM & 0.0377 & 0.0614 & 0.0244 & 0.0320 & 0.0536 & 0.0921 & 0.0327 & 0.0450 & 0.0069 & 0.0131 & 0.0043 & 0.0063 \\
SEMCo & 0.1631 & 0.2306 & 0.1198 & 0.1416 & 0.0907 & 0.1562 & 0.0563 & 0.0773 & 0.0204 & 0.0354 & 0.0126 & 0.0174 \\
\midrule
\textbf{CKG-RL} & \textbf{0.2473}\sigmark{*} & \textbf{0.2863}\sigmark{*} & \textbf{0.1972}\sigmark{*} & \textbf{0.2098}\sigmark{*} & \textbf{0.1724}\sigmark{*} & \textbf{0.2443}\sigmark{*} & \textbf{0.1137}\sigmark{*} & \textbf{0.1370}\sigmark{*} & \textbf{0.0487}\sigmark{*} & \textbf{0.0679} & \textbf{0.0330}\sigmark{*} & \textbf{0.0392}\sigmark{*} \\
\emph{Imp.} & +16.6\% & +10.6\% & +16.3\% & +13.7\% & +29.5\% & +16.2\% & +32.4\% & +23.9\% & +23.9\% & +3.5\% & +33.1\% & +18.1\% \\
\bottomrule
\end{tabular}
}
\caption{Strict course-cold performance (three-seed means). Underlines mark second best. \(^{*}\): starred cells pass a Holm-corrected paired Wilcoxon test against the strongest displayed reference (MOOCCube/Junyi/COCO: CGRC/ALDI/CCFCRec), all with \(p_{\mathrm{Holm}}<0.01\); COCO R@10 is unmarked. Imp.: gain over the same reference.}
\label{tab:main-cold-results}
\end{table*}

Baselines are ordered chronologically so that the most recent methods are closest to CKG-RL. Popularity remains first as a leakage check and receives zero cold course-macro score because held-out courses have no training popularity. BPR and LightGCN recover learner-side signal but lack target-course edges. Content-transfer and graph-reconstruction baselines borrow side information or warm-course structure. CKG-RL adds protocol-aligned masking plus course knowledge in simulation and rewards. For statistical evidence, Holm-corrected tests reject the matched-pair null for all four MOOCCube @5/@10 metrics against CGRC and all four Junyi metrics against ALDI. On COCO, the tests reject for R@5, N@5, and N@10 against CCFCRec; R@10 is a descriptive mean gain without a significance mark. The supplement reports seed-level CIs, all paired \(p\)-values, and the more conservative all-metric correction.

\subsection{RQ2: Exposure and Pedagogical-Risk Diagnostics}
Figure~\ref{fig:method-motivation} is a diagnosis of existing methods rather than a second performance comparison. On the held-out Top-10 lists, PCGNN has mean prerequisite and difficulty gaps of 0.6427 and 0.0963, respectively, showing that course-specific modeling alone does not remove structural mismatch for zero-interaction targets. The CGRC course-level distribution exposes a different failure mode: 46\% of the 204 audited seed-course cases have NDCG@10 no greater than 0.10, and the mean cold-course share in its Top-10 lists is 24.6\%. CKG-RL responds to the first limitation with knowledge-guided sampling, educational rewards, and prerequisite supervision, and to the second with content anchoring, cold-ID masking, and learner simulation. Figure~\ref{fig:p1-topk-motivation} then tests which outcomes improve and which trade-offs remain.

For the complete post-hoc audit, we score the actual Top-10 lists emitted by the frozen CKG-RL, PCGNN, and CGRC checkpoints. Each recommended course is evaluated from training history only by prerequisite coverage gap, directional concept overlap, a content-structural difficulty gap, and structural redundancy. Structural complexity equally weights robustly normalized prerequisite and concept counts; readiness averages up to five available, structurally most advanced historical courses. Redundancy is the maximum same-family duplication or bidirectional video-unit containment. We average over ranks, then interactions with the same cold target course, and compare CKG-RL separately with each baseline on 204 matched (seed, target-course) units. Lists without a recommended cold course remain missing in the cold-only analysis.

\begin{figure}[!t]
\centering
\includegraphics[width=\columnwidth]{figures/mooccube_p1_topk_motivation.pdf}
\caption{Complete Top-10 exposure and model-neutral pedagogical-risk audit on MOOCCube. (a) Circles compare CKG-RL with course-specific PCGNN and squares compare it with generic cold-start CGRC; positive paired effects favor CKG-RL after orienting lower-is-better risks. (b) Cold-course proportion for all three models, shown as three-seed mean and standard deviation.}
\label{fig:p1-topk-motivation}
\end{figure}

Figure~\ref{fig:p1-topk-motivation} exposes distinct trade-offs against the two baseline classes. Against PCGNN, CKG-RL lowers prerequisite and difficulty gaps and increases cold-course share by 8.5 percentage points, but its concept-continuity difference is inconclusive and structural redundancy is higher by 0.0117. Against CGRC, CKG-RL lowers prerequisite gap and structural redundancy and increases cold-course share by 26.2 percentage points, but concept continuity is lower by 0.0317 and the difficulty difference is inconclusive. The full audit therefore supports a joint motivation built from complementary baseline weaknesses, not a claim that CKG-RL dominates either reference on every pedagogical proxy.

\subsection{RQ3: Course-Knowledge Ablation}
\begin{table*}[!t]
\centering
{
\footnotesize
\setlength{\tabcolsep}{5.5pt}
\begin{tabular*}{\textwidth}{@{\extracolsep{\fill}}lcccccc}
\toprule
& \multicolumn{3}{c}{Recall \(\uparrow\)} & \multicolumn{3}{c}{NDCG \(\uparrow\)} \\
\cmidrule(lr){2-4}\cmidrule(lr){5-7}
Variant & @5 & @10 & @20 & @5 & @10 & @20 \\
\midrule
\textbf{Full CKG-RL} & \textbf{0.2473} & \textbf{0.2863} & \textbf{0.3220} & \textbf{0.1972} & \textbf{0.2098} & \textbf{0.2188} \\
w/o educ. rewards & 0.2409 & 0.2791 & 0.3159 & 0.1929 & 0.2054 & 0.2146 \\
w/o KG sampler & 0.2430 & 0.2810 & 0.3201 & 0.1931 & 0.2055 & 0.2154 \\
w/o prereq. aux. loss & 0.2422 & 0.2779 & 0.3160 & 0.1911 & 0.2028 & 0.2124 \\
w/o course-KG inputs & 0.2451 & 0.2804 & 0.3161 & 0.1937 & 0.2053 & 0.2142 \\
\bottomrule
\end{tabular*}
}
\caption{Component ablation on MOOCCube under strict course-cold evaluation (three-seed means). Best values are bold.}
\label{tab:core-ablation}
\end{table*}

Table~\ref{tab:core-ablation} reports core ablations on MOOCCube using the same split, candidate set, \(K\) values, and course-macro aggregation as the main comparison. Removing educational rewards reduces R@10/N@10 by 0.0072/0.0044. Removing the knowledge-guided sampler reduces them by 0.0053/0.0043, and removing the prerequisite auxiliary loss gives the largest early-rank NDCG drop, decreasing N@5/N@10 by 0.0061/0.0070. When all course-knowledge inputs are removed, scores drop below Full CKG-RL on all six metrics. This pattern points to course knowledge coupled with rewards, candidate construction, and auxiliary supervision rather than a gain from one scalar signal.

We further compare the frozen Full CKG-RL lists with matched \emph{w/o educational rewards} and \emph{w/o simulator} exports using the same model-neutral risk audit as RQ2. Over 204 paired seed-course units, the educational reward has a targeted rather than uniform effect: among recommended cold courses, Full CKG-RL lowers prerequisite gap by 0.0132 (95\% CI [0.0056, 0.0205] after orienting improvement as positive) and difficulty gap by 0.0034 ([0.0011, 0.0056]), while increasing cold-course exposure by 2.7 percentage points. Its overall prerequisite and difficulty effects are inconclusive, and concept continuity decreases. The simulator comparison provides no positive risk effect under these proxies; continuity and redundancy move adversely in parts of the audit. In the same export replay, removing the simulator lowers the three-seed course-macro R@10/N@10 means from 0.2863/0.2098 to 0.2839/0.2078. We therefore use only the supported course-reward effects as pedagogical mechanism evidence; the simulator remains an accuracy/exposure component whose pedagogical role is not established by this audit.

The supplementary motivation-evidence table consolidates only the predeclared claims supported by non-crossing paired bootstrap intervals. Cold-course share is reported as an exposure diagnostic rather than a learning-quality outcome; inconclusive and adverse results remain in the complete audit above and are not used as positive motivation claims.

\subsection{RQ4: Sensitivity and Cost}
\begin{figure}[!t]
\centering
\includegraphics[width=\columnwidth]{figures/mooccube_hparam_validation_updated_default_aaai_main.pdf}
\caption{Hyperparameter sensitivity on MOOCCube. Panels report validation cold course-macro Recall@10 and NDCG@10; dotted vertical lines denote the reported default \((\beta=0.20,\) reward scale \(=1.00,\ T=5)\).}
\label{fig:hparam-sensitivity}
\end{figure}

The MOOCCube sweeps vary knowledge-sampling weight \(\beta\), reward-shaping scale, and simulator horizon \(T\) on validation folds around the reported default \((\beta=0.20,\) reward scale \(=1.00,\ T=5)\). The default gives validation R@10 \(=0.2942\) and N@10 \(=0.2175\). Across \(\beta\), validation R@10 spans 0.2831--0.2974 and N@10 spans 0.2117--0.2194. Reward scale stays narrow, with R@10/N@10 ranging over 0.2864--0.2944/0.2126--0.2182. Horizons \(T=1,3,5,7,10\) also remain close; \(T=7\) gives the lowest validation N@10 at 0.2105. These validation sweeps are reported as sensitivity evidence around the fixed main-table configuration, and Table~\ref{tab:main-cold-results} reports held-out test scores without test-set hyperparameter selection.

Table~\ref{tab:efficiency-aaai} reports wall-clock training and inference cost against CGRC, the closest graph-reconstruction cold-start baseline. CKG-RL averages \(95.6\pm21.9\), \(312.2\pm67.7\), and \(279.2\pm17.8\) seconds per epoch on MOOCCube, Junyi, and COCO. Eval-only final cold+hot inference takes \(18.0\pm0.7\), \(56.4\pm0.3\), and \(9.8\pm0.1\) seconds, respectively. The displayed Junyi CGRC train/epoch value is a provisional cross-dataset projection pending the prioritized three-seed profile. The MOOCCube inference profile uses the main split with a retained checkpoint that has the same CKG-RL inference graph. The simulator mainly affects training; cached inference remains \(O(QId)\).

\begin{table*}[t]
\centering
{
\small
\setlength{\tabcolsep}{2.4pt}
\renewcommand{\arraystretch}{1.02}
\begin{tabular*}{\textwidth}{@{\extracolsep{\fill}}llccccc}
\toprule
Dataset & Cost ref. & CKG train/epoch & CGRC train/epoch & CKG infer. & CGRC infer. & \(\Delta\)N@10 vs CGRC [95\% CI] \\
\midrule
MOOCCube & CGRC & 95.6\,$\pm$\,21.9 & 70.5\,$\pm$\,4.6 & 18.0\,$\pm$\,0.7 & 22.2\,$\pm$\,0.4 & +0.0263 [0.0114, 0.0401] \\
Junyi & CGRC & 312.2\,$\pm$\,67.7 & 333.0\,$\pm$\,45.0 & 56.4\,$\pm$\,0.3 & 69.2\,$\pm$\,0.2 & +0.0330 [0.0159, 0.0512] \\
COCO & CGRC & 279.2\,$\pm$\,17.8 & 389.7\,$\pm$\,80.0 & 9.8\,$\pm$\,0.1 & 11.8\,$\pm$\,0.2 & +0.0183 [0.0162, 0.0205] \\
\bottomrule
\end{tabular*}
}
\caption{Efficiency comparison with CGRC.}
\label{tab:efficiency-aaai}
\vspace{2pt}
\parbox{\textwidth}{\footnotesize Times are mean\(\pm\)std seconds over retained seed-level profiles. \(\Delta\)N@10 is the mean paired course-level difference with a 95\% bootstrap CI over matched seed-course units pooled across three seeds; intervals are conditional on the three trained model fits.}
\end{table*}


\section{Discussion and Limitations}
The experiments point to the evaluation unit as a source of hidden risk. Interaction-weighted metrics can improve when a model serves a few high-volume cold courses even if many others receive no useful rank~\cite{ekstrand2024distributional,bauer2024landscape}; course-macro full ranking exposes this pattern. CGRC already favors continuity and smaller difficulty gaps, but its prerequisite alignment is flat and redundancy is concentrated near the top. Course knowledge is therefore useful as a constraint beyond content similarity.

The recommendation-list audit also separates exposure from pedagogical fit. CKG-RL roughly doubles the share of cold courses in Top-10 lists and reduces prerequisite gap and early-rank redundancy, but this broader exposure coincides with lower concept continuity and a cold-only difficulty disadvantage. We therefore do not combine the four proxies into a post-hoc scalar risk score: no composite weights were preregistered, and such a score could hide effects with opposite directions. These results support a claim about exposure and specific structural signals, not a claim that CKG-RL uniformly improves pedagogical appropriateness.

The size of the gain tracks the quality of the available course knowledge. Improvements are largest on MOOCCube and Junyi, where MOOCCubeX supplies human-curated prerequisite and concept relations, and smallest on COCO, whose relations are derived from course metadata rather than expert annotation and whose roughly 6{,}000 warm candidates per seed dilute each cold course across the largest catalog. The educational reward therefore contributes most when the underlying prerequisite and concept signals are reliable, which is consistent with our central claim rather than a counterexample: CKG-RL converts course knowledge into ranking behavior, so its benefit scales with how informative that knowledge is.

Deployment validation is outside the scope of this offline study. The benchmarks isolate about 68 MOOCCube cold courses, 71 Junyi cold courses, and 820 COCO cold courses per seed, a smaller scale than industrial catalogs. Course relations and embeddings are dataset-specific; all four list-level signals are proxies rather than direct mastery or learning-outcome measurements. The 204 course-level pairs share only three trained model fits, so their uncertainty is conditional on those fits rather than evidence of broad seed-level generalization. CGRC is a same-configuration repaired rerun rather than a bitwise reproduction of the old main-table checkpoint; its seed-2026 training resumed without RNG-state restoration, and checkpoint replay changes interaction NDCG@10 by at most $2.26\times10^{-4}$. Deployment therefore requires online and pedagogical validation~\cite{dasilva2023educationalrs,stray2024values}.

\section{Conclusion}
We studied strict course-cold MOOC recommendation, where target courses have no training interactions and full-ranking course-macro evaluation exposes failures hidden by interaction averages. CKG-RL combines content anchoring, forced-cold masking, retrieval-based simulation, and course-knowledge rewards. Across MOOCCube, Junyi, and COCO, it yields higher mean early-rank course-cold metrics under the same full-ranking course-macro protocol. On actual MOOCCube Top-10 lists, it roughly doubles cold-course exposure and improves prerequisite coverage and early-rank redundancy, but sacrifices concept continuity and does not uniformly reduce pedagogical risk. This trade-off defines a concrete target for future constrained or multi-objective course ranking.

\section*{Ethical Statement}
Experiments use public MOOCCube/MOOCCubeX, Junyi, and COCO data plus derived course representations. Released artifacts should exclude personally identifiable information and follow dataset terms. Deployment requires online validation and monitoring for visibility, fairness, transparency, and consequential use~\cite{stray2024values,zhao2025fairnessdiversity}.

\bibliography{references}

\end{document}
